% Propietario conservado: /Users/ruben/Documents/New project/output/EXTREMA_MEDIA_RAZON_RECIPROCIDAD_ES_EN_20260918/ARTICULO/ES/source/libro/certificacion_lean.tex
\chapter*{Formalización aritmética del balance bilateral}
\addcontentsline{toc}{section}{Formalización aritmética del balance bilateral}
\label{x_reciprocidad:libro:formalizacion-lean}

La especialización nonádica de la ley bilateral admite una formalización
aritmética independiente. El punto de aplicación se encuentra después de la
construcción de sus coeficientes y transportes: la partición de base diez fija
\(a=8\), y los dos registros son \(8x-y\) y \(9x+y\). Sobre ese dominio,
Lean~4.21.0 verifica siete teoremas con su biblioteca \texttt{Std}. Las fuentes
formales y el procedimiento de ejecución se proporcionan como suplemento.
El presente complemento enuncia su contenido matemático y lo relaciona con
las demostraciones operatorias del artículo.

\section*{Identidad polinómica y conservación de la norma}

Para cualesquiera \(x,y\in\mathbb Z\), se cumple
\[
 9(8x-y)^2+8(9x+y)^2=17(72x^2+y^2).
\]
La expansión de los dos cuadrados cancela los productos cruzados. Cuando
\(y^2=x^2\), la misma igualdad da
\[
 9(8x-y)^2+8(9x+y)^2=17\cdot73\,x^2.
\]
Éstos son los dos primeros teoremas formalizados. El primero expresa un
balance para entradas enteras arbitrarias; el segundo aplica la hipótesis
explícita de conservación de la norma. El entero \(73=72+1\) ocupa
exactamente el lugar que tiene en la \cref{x_reciprocidad:lib04:teo-balance}.

\section*{Dimensión finita arbitraria}

Sean \(n\in\mathbb N\) y \(x,y:\mathbb N\to\mathbb Z\). Definamos
\[
 S_n(x)=\sum_{j<n}x_j^2,
 \qquad
 E_n(x,y)=\sum_{j<n}\bigl(9(8x_j-y_j)^2+8(9x_j+y_j)^2\bigr).
\]
Los dos teoremas siguientes establecen
\[
 E_n(x,y)=17\bigl(72S_n(x)+S_n(y)\bigr)
\]
y, bajo la hipótesis \(S_n(y)=S_n(x)\),
\[
 E_n(x,y)=17\cdot73\,S_n(x).
\]
La prueba formal procede por inducción sobre \(n\). El caso inicial es la
suma vacía; el paso de inducción añade una coordenada y aplica la identidad
polinómica ya demostrada. El cuantificador universal sobre \(n\) incluye
todas las dimensiones finitas. La hipótesis utiliza la suma total de
cuadrados, por lo que admite redistribuciones entre coordenadas.

\section*{Inversión exacta de los canales bilaterales}

Escribamos \(q_-=8x-y\) y \(q_+=9x+y\). Los teoremas quinto y sexto
son las igualdades exactas
\[
 q_-+q_+=17x,\qquad 8q_+-9q_-=17y.
\]
En la imagen de la aplicación \((x,y)\mapsto(q_-,q_+)\), ambos miembros
izquierdos son múltiplos de \(17\). Las divisiones recuperan las dos
entradas. El séptimo teorema formula su consecuencia inyectiva:
\[
 \left.
 \begin{aligned}
 8x-y&=8x'-y',\\
 9x+y&=9x'+y'
 \end{aligned}
 \right\}
 \quad\Longrightarrow\quad x=x'\ \text{y}\ y=y'.
\]
La cancelación en \(\mathbb Z\) justifica esta conclusión. El contenido
formal coincide con la recuperación del par de registros, antes de
normalizarlos como canales de una realización métrica.

\section*{Correspondencia entre demostración algebraica y formalización}

La formalización utiliza variables enteras arbitrarias e inducción
matemática. Por ello constituye una prueba de los enunciados anteriores,
distinta de la comprobación de un conjunto finito de ejemplos. El sistema
acepta los términos de prueba sin declaraciones axiomáticas añadidas ni
demostraciones diferidas. El suplemento registra asimismo las dependencias
lógicas comunicadas por la herramienta para cada declaración.

La ley operatoria del cuerpo del artículo posee un dominio más amplio:
isometrías entre espacios de Hilbert, archivo de profundidad arbitraria y
lectores conjugados. Sus demostraciones se encuentran en las
\cref{x_reciprocidad:lib04:teo-balance,x_reciprocidad:lib06:simetria,x_reciprocidad:lib06:profundidad}. La prueba Lean
aquí incorporada certifica la especialización aritmética enunciada, con
sus cuantificadores y su estructura de partida. La procedencia de los
coeficientes, la generación del registro y las realizaciones excepcionales
conservan las demostraciones anteriores de las que dependen.
